\subsection{定义在集族上的测度的延拓}

设 \(\Phi\) 是局部紧空间 \(X\) 的子集所成的一个非空集。给定一个实值函数 \(M \mapsto \alpha(M)\)，它有定义且 \(\geqslant 0\) 于 \(\Phi\) 上，我们提出寻求这样的条件：在这些条件下，存在一个正测度 \(\mu\)（\(X\) 上的），使得 \(\Phi\) 中的集合都是 \(\mu\)-可积的，并且 \(\mu(M) = \alpha(M)\) 对所有 \(M \in \Phi\) 成立。我们将只考虑集 \(\Phi\) 满足下列条件的情形：

\((\mathrm{PC}_{\mathrm{I}})\) \(\Phi\) 中两个集合的并与交属于 \(\Phi\)。

\((\mathrm{PC}_{\mathrm{II}})\) 对每个由紧集 \(K\) 与开集 \(U\)（\(X\) 中满足 \(K \subset U\) 者）组成的对，存在一个集合 \(M \in \Phi\)，使得 \(K \subset M \subset U\)。

注意，条件 \((\mathrm{PC}_{\mathrm{II}})\) 蕴涵 \(\varnothing\in\Phi\)，这只需取 \(K=U=\varnothing\)。然而，集 \(\Phi\) 未必是一个集环；例如，X 的所有紧子集所成之集满足条件 \((\mathrm{PC}_{\mathrm{I}})\) 与 \((\mathrm{PC}_{\mathrm{II}})\)，但一般说来不是一个集环，因为若 M 与 N 是紧的，则 \(M\cap \complement N\) 一般说来并不如此。

我们还将假设函数 \(\alpha\)（定义在 \(\Phi\) 上）满足下列条件（它们显然是该问题有解的必要条件）：

\((\mathrm{PM}_{\mathrm{I}})\) 关系 \(\mathrm{M}\subset \mathrm{N}\) 蕴涵 \(\alpha (\mathrm{M})\leqslant \alpha (\mathrm{N})\)

\((\mathrm{PM}_{\mathrm{II}})\) 对 \(\Phi\) 中任意的 M 与 N，\(\alpha (\mathrm{M}\cup \mathrm{N})\leqslant \alpha (\mathrm{M}) + \alpha (\mathrm{N})\)

\((\mathrm{PM}_{\mathrm{III}})\) 关系 \(\mathrm{M}\cap \mathrm{N} = \emptyset\) 蕴涵 \(\alpha (\mathrm{M}\cup \mathrm{N}) = \alpha (\mathrm{M}) + \alpha (\mathrm{N})\)

取 \(\mathbf{N} = \varnothing\) 于条件 \((\mathrm{PM}_{\mathrm{III}})\) 中，可推出 \(\alpha(\varnothing) = 0\)；于是条件 \((\mathrm{PM}_{\mathrm{I}})\) 表明，\(\alpha(M) \geqslant 0\) 对每个 \(M \in \Phi\) 成立。

定理 5. --- 设 \(\Phi\) 是局部紧空间 \(X\) 的子集所成的一个集，满足 \((\mathrm{PC}_{\mathrm{I}})\) 与 \((\mathrm{PC}_{\mathrm{II}})\)，并设 \(\alpha\) 是定义在 \(\Phi\) 上的实值函数，满足条件 \((\mathrm{PM}_{\mathrm{I}})\)、\((\mathrm{PM}_{\mathrm{II}})\) 与 \((\mathrm{PM}_{\mathrm{III}})\)。为使存在一个正测度 \(\mu\)（\(X\) 上的），使得 \(\Phi\) 中的集合都是 \(\mu\)-可积的，并且 \(\mu(M) = \alpha(M)\) 对所有 \(M \in \Phi\) 成立，其充要条件是 \(\alpha\) 还满足下列条件：

\((PM_{IV})\) 对每个 \(\varepsilon > 0\) 与每个 \(M \in \Phi\)，存在一个紧集 \(K \subset M\) 与一个开集 \(U \supset M\)，使得对每个 \(N \in \Phi\)（满足关系 \(K \subset N \subset U\) 者），有 \(|\alpha(N) - \alpha(M)| \leqslant \varepsilon\)。

此外，若条件 \((\mathrm{PM}_{\mathrm{IV}})\) 得到满足，则测度 \(\mu\) 是唯一的；对每个紧集 K，\(\mu(\mathrm{K}) = \inf_{\mathrm{M} \in \Phi, \mathrm{M} \supset \mathrm{K}} \alpha(\mathrm{M})\)；对每个开集 U，\(\mu^{*}(\mathrm{U}) =\)

sup \(\alpha (\mathbf{M})\)

\(\mathbf{M}\in \Phi ,\bar{\mathbf{M}}\subset \mathbf{U}\)

注意，条件 \((PM_{IV})\) 等价于下列两个条件的合取：

\((\mathrm{PM}_{\mathrm{IV}}^{\prime})\) 对每个 \(\varepsilon > 0\) 与每个 \(\mathbf{M} \in \Phi\)，存在一个开集 \(\mathbf{U} \supset \mathbf{M}\)，使得对每个 \(\mathbf{N} \in \Phi\)（含于 \(\mathbf{U}\) 者），有 \(\alpha(\mathbf{N}) \leqslant \alpha(\mathbf{M}) + \varepsilon\)。 \((\mathrm{PM}_{\mathrm{IV}}^{\prime \prime})\) 对每个 \(\varepsilon > 0\) 与每个 \(\mathbf{M} \in \Phi\)，存在一个紧集 \(\mathbf{K} \subset \mathbf{M}\)，使得对每个 \(\mathbf{N} \in \Phi\)（包含 \(\mathbf{K}\) 者），有 \(\alpha(\mathbf{N}) \geqslant \alpha(\mathbf{M}) - \varepsilon\)。

因为，显然 \((\mathrm{PM}_{\mathrm{IV}}^{\prime})\) 与 \((\mathrm{PM}_{\mathrm{IV}}^{\prime\prime})\) 蕴涵 \((\mathrm{PM}_{\mathrm{IV}})\)。反之，例如让我们证明 \((\mathrm{PM}_{\mathrm{IV}})\) 蕴涵 \((\mathrm{PM}_{\mathrm{IV}}^{\prime})\)：设 K 是一个紧集，U 是一个开集，使得 \(K \subset M \subset U\)，并且 \(|\alpha(P) - \alpha(M)| \leqslant \varepsilon\) 对每个 \(P \in \Phi\)（满足 \(K \subset P \subset U\) 者）成立。那么，若 \(N \in \Phi\) 且 \(N \subset U\)，则 \(M \cup N\) 属于 \(\Phi\) 且 \(K \subset M \cup N \subset U\)，于是 \(\alpha(M \cup N) \leqslant \alpha(M) + \varepsilon\)，从而更有 \(\alpha(N) \leqslant \alpha(M) + \varepsilon\)。

当集 \(\Phi\)（满足 \((\mathrm{PC}_{1})\) 与 \((\mathrm{PC}_{11})\) 的）由紧集组成时，条件 \((\mathrm{PM}_{\mathrm{IV}}^{\prime \prime})\) 平凡地成立，于是 \((\mathrm{PM}_{\mathrm{IV}})\) 等价于 \((\mathrm{PM}_{\mathrm{IV}}^{\prime})\)。

条件 \((\mathrm{PM}_{\mathrm{IV}})\) 是必要的：这由 No.~6 的定理 4（关于用紧集与开集“逼近”可积集）立即得出。为证明定理的其余断言，我们分几步进行。

\(1^{\circ}\) \(\mathfrak{P}(\mathrm{X})\) 上一个拓扑的定义。

对由 X 的紧集 K 与开集 U 组成的每个对 \((K,U)\)，我们用 \(I(K,U)\) 表示子集 \(M \subset X\)（满足 \(K \subset M \subset U\) 者）所成之集；为使 \(I(K,U)\) 非空，其充要条件是 \(K \subset U\)。若 \((K',U')\) 是由紧集 \(K'\) 与开集 \(U'\) 组成的第二个对，则我们有

\[
\mathrm{I} (\mathrm{K}, \mathrm{U}) \cap \mathrm{I} (\mathrm{K} ^ {\prime}, \mathrm{U} ^ {\prime}) = \mathrm{I} (\mathrm{K} \cup \mathrm{K} ^ {\prime}, \mathrm{U} \cap \mathrm{U} ^ {\prime}).
\]

设 \(\mathcal{T}\) 是 \(\mathfrak{P}(\mathrm{X})\) 上由诸子集 I(K,U)（当 K 取遍 X 的紧子集所成之集而 U 取遍开子集所成之集时）生成的拓扑；由前所述，诸 I(K,U) 构成拓扑 \(\mathcal{T}\) 的一个基（GT, I, §1, No.~3）。

我们注意到，\(\mathcal{T}\) 的定义蕴涵：在 \(\mathfrak{P}(\mathrm{X})\) 中，\(\mathrm{X}\) 的紧子集所成之集是稠密的。条件 \((\mathrm{PC}_{\mathrm{II}})\) 表达的是 \(\Phi\) 在 \(\mathfrak{P}(\mathrm{X})\) 中稠密，而条件 \((\mathrm{PM}_{\mathrm{IV}})\) 表达的是函数 \(\alpha\) 在 \(\Phi\) 上关于由 \(\mathcal{T}\) 诱导的拓扑连续。最后，No.~6 的定理 4 表达的是：函数 \(\mathrm{M} \mapsto \mu(\mathrm{M})\) 在 \(\mu\)-可积集所成的集环上关于由 \(\mathcal{T}\) 诱导的拓扑连续。

\(2^{\circ}\) \(\mu\) 的唯一性。

我们用 \(\overline{\Phi}\) 表示满足下述条件的子集 \(M \subset X\) 所成之集：\(\alpha(N)\) 当 N（关于拓扑 T）趋于 M 而保持在 \(\Phi\) 中时，趋于一个有限极限；于是我们可以把 \(\alpha\) 以唯一的方式延拓为一个连续映射 \(\overline{\alpha}\)（\(\overline{\Phi}\) 到 R 中的）（GT, I, §8, No.~5, 定理 1）。若存在一个满足所要求条件的测度 \(\mu\)，则上面的注记证明，集环 \(\Psi\)（\(\mu\)-可积集所成者）含于 \(\overline{\Phi}\)，并且 \(\mu(M) = \overline{\alpha}(M)\) 对每个 \(M \in \Psi\) 成立；这一关系特别地对 X 的每个紧子集 M 成立，这就证明了 \(\mu\) 的唯一性（No.~10, 命题 19 的推论 3）。

\(3^{\circ}\) \(\alpha\) 到紧集的延拓。

现在我们将在不假设 \(\mu\) 存在的情况下，研究集 \(\overline{\Phi}\) 以及延拓 \(\overline{\alpha}\)（\(\alpha\) 到 \(\overline{\Phi}\) 上的）。我们首先证明每个紧集 \(K\) 都属于 \(\overline{\Phi}\)，并且 \(\overline{\alpha}(K) = \inf_{P \in \Phi, P \supset K} \alpha(P)\)。设 \(a = \inf_{P \in \Phi, P \supset K} \alpha(P)\)；对每个 \(\varepsilon > 0\)，存在一个 \(M \in \Phi\)，使得 \(K \subset M\) 且 \(\alpha(M) \leqslant a + \varepsilon\)。由 \((PM_{IV}')\)，存在一个开集 \(U \supset M\)，使得对每个 \(N \in \Phi\)（含于 \(U\) 者），有 \(\alpha(N) \leqslant \alpha(M) + \varepsilon \leqslant a + 2\varepsilon\)；于是对每个 \(N \in \Phi\)（满足 \(K \subset N \subset U\) 者），有 \(a \leqslant \alpha(N) \leqslant a + 2\varepsilon\)，这按定义表明 \(K \in \overline{\Phi}\) 且 \(\overline{\alpha}(K) = a\)。

这一结果立即表明，若 \(K_{1}\) 与 \(K_{2}\) 是使 \(K_{1} \subset K_{2}\) 的两个紧集，则 \(\overline{\alpha}(K_{1}) \leqslant \overline{\alpha}(K_{2})\)。若 \(K_{1}\) 与 \(K_{2}\) 是任意两个紧集，则有 \(\overline{\alpha}(K_{1} \cup K_{2}) \leqslant \overline{\alpha}(K_{1}) + \overline{\alpha}(K_{2})\)（由 \((\mathrm{PM}_{\mathrm{II}})\)）。此外，我们还将看到，若 \(K_{1}\) 与 \(K_{2}\) 不相交，则 \(\overline{\alpha}(K_{1} \cup K_{2}) = \overline{\alpha}(K_{1}) + \overline{\alpha}(K_{2})\)。因为此时存在两个不相交的开集 \(U_{1}, U_{2}\)，使 \(K_{1} \subset U_{1}, K_{2} \subset U_{2}\)（GT, II, §4, 命题 4）。于是由 \((\mathrm{PC}_{\mathrm{II}})\)，还存在两个集合 \(M_{1} \in \Phi, M_{2} \in \Phi\)，使 \(K_{1} \subset M_{1} \subset U_{1}\) 且 \(K_{2} \subset M_{2} \subset U_{2}\)。现设 P 是 \(\Phi\) 中包含 \(K_{1} \cup K_{2}\) 的任一集合；两个集合 \(P \cap M_{1}\) 与 \(P \cap M_{2}\) 的并属于 \(\Phi\)（由 \((\mathrm{PC}_{\mathrm{I}})\)），且由于这两个集合不相交，应用 \((\mathrm{PM}_{\mathrm{I}})\) 与 \((\mathrm{PM}_{\mathrm{III}})\) 便得

\[
\alpha (\mathrm{P}) \geqslant \alpha (\mathrm{P} \cap \mathrm{M} _ {1}) + \alpha (\mathrm{P} \cap \mathrm{M} _ {2}) \geqslant \overline {{\alpha}} (\mathrm{K} _ {1}) + \overline {{\alpha}} (\mathrm{K} _ {2}),
\]

由此便建立了我们的断言。

\(4^{\circ}\) \(\alpha\) 到开集的延拓。

现在我们将看到，为使一个开集 U 属于 \(\overline{\Phi}\)，其充要条件是：当 K 取遍 U 的紧子集所成之集时，诸数 \(\overline{\alpha}(K)\) 的上确界有限；此外，此时 \(\overline{\alpha}(U)\) 等于这个上确界。

因为，设 U 是属于 \(\bar{\Phi}\) 的一个开集；对每个 \(\varepsilon > 0\)，存在一个紧集 \(K \subset U\)，使得对每个集合 \(M \in \Phi\)（满足 \(K \subset M \subset U\) 者），有 \(|\overline{\alpha}(U) - \alpha(M)| \leqslant \varepsilon\)，从而 \(|\overline{\alpha}(U) - \overline{\alpha}(K)| \leqslant \varepsilon\)；另一方面，若 \(K'\) 是含于 U 的任一紧集，则 \(K \subset K \cup K' \subset U\)，于是 \(|\overline{\alpha}(U) - \overline{\alpha}(K \cup K')| \leqslant \varepsilon\)，从而 \(\overline{\alpha}(U) \geqslant \overline{\alpha}(K \cup K') - \varepsilon \geqslant \overline{\alpha}(K') - \varepsilon\)；因此 \(\overline{\alpha}(U)\) 确实等于当 K 取遍 U 的紧子集所成之集时诸数 \(\overline{\alpha}(K)\) 的上确界。

反之，设 U 是使 \(b = \sup_{K \subset U} \overline{\alpha}(K) < +\infty\) 的一个开集（其中 K 取

遍 U 的紧子集所成之集），我们来证明 \(U \in \overline{\Phi}\)。对每个 \(\varepsilon > 0\)，存在一个紧集 \(K \subset U\)，使 \(b - \varepsilon \leqslant \overline{\alpha}(K) \leqslant b\)；由 \((\mathrm{PM}_{\mathrm{IV}}^{\prime\prime})\)，对每个集合 \(M \in \Phi\)（满足 \(K \subset M \subset U\) 者），存在一个紧集 \(K' \subset M\)，使得

\[
\alpha (\mathrm{M}) \leqslant \overline {{{{\alpha}}}} (\mathrm{K} ^ {\prime}) + \varepsilon \leqslant b + \varepsilon ;
\]

因此 \(b - \varepsilon \leqslant \alpha (\mathrm{M}) \leqslant b + \varepsilon\)，这就证明了 \(\mathrm{U} \in \overline{\Phi}\)。

由开集 \(U \in \overline{\Phi}\) 以及 \(\overline{\alpha}(U)\) 的这一刻画，首先得出：若 \(U_{1}\) 与 \(U_{2}\) 是使 \(U_{1} \subset U_{2}\) 且 \(U_{2} \in \overline{\Phi}\) 的两个开集，则 \(U_{1} \in \overline{\Phi}\) 且 \(\overline{\alpha}(U_{1}) \leqslant \overline{\alpha}(U_{2})\)。另一方面，若 \(U_{1}\) 与 \(U_{2}\) 是属于 \(\overline{\Phi}\) 的两个开集，则 \(U_{1} \cup U_{2}\) 亦然，且 \(\overline{\alpha}(U_{1} \cup U_{2}) \leqslant \overline{\alpha}(U_{1}) + \overline{\alpha}(U_{2})\)。因为，设 K 是含于 \(U_{1} \cup U_{2}\) 的任一紧集；对每个点 \(x \in K\)，存在 x 的一个含于 \(U_{1}\) 或 \(U_{2}\) 之一的紧邻域；因此可以用有限个这样的邻域覆盖 K；若 \(K_{1}\) （相应地 \(K_{2}\)）是含于 \(U_{1}\) （相应地 \(U_{2}\)）的那些邻域的并，则 \(K \subset K_{1} \cup K_{2}\)，于是

\[
\overline {{{{\alpha}}}} (\mathrm{K}) \leqslant \overline {{{{\alpha}}}} (\mathrm{K} _ {1} \cup \mathrm{K} _ {2}) \leqslant \overline {{{{\alpha}}}} (\mathrm{K} _ {1}) + \overline {{{{\alpha}}}} (\mathrm{K} _ {2}) \leqslant \overline {{{{\alpha}}}} (\mathrm{U} _ {1}) + \overline {{{{\alpha}}}} (\mathrm{U} _ {2}),
\]

由此便建立了所述的性质。

\(5^{\circ}\) \(\overline{\Phi}\) 与 \(\overline{\alpha}\) 的性质。

\(\overline{\Phi}\) 与 \(\overline{\alpha}\) 的定义现在可以改述如下（考虑到 \((\mathrm{PC}_{\mathrm{II}})\)）：为使 \(M \in \overline{\Phi}\)，其充要条件是，对每个 \(\varepsilon > 0\)，存在一个紧集 K 与一个开集 \(U \in \overline{\Phi}\)，使得 \(K \subset M \subset U\) 且 \(\overline{\alpha}(U) - \overline{\alpha}(K) \leqslant \varepsilon\)；此外，\(\overline{\alpha}(M)\) 是诸 \(\overline{\alpha}(U)\)（对应于包含 M 的开集 \(U \in \overline{\Phi}\)）的下确界，以及诸 \(\overline{\alpha}(K)\)（对应于紧集 \(K \subset M\)）的上确界。

由此我们首先推出：若 \(\mathrm{M}_1\)、\(\mathrm{M}_2\) 与 \(\mathrm{M}_1 \cup \mathrm{M}_2\) 属于 \(\overline{\Phi}\)，则 \(\overline{\alpha}(\mathrm{M}_1 \cup \mathrm{M}_2) \leqslant \overline{\alpha}(\mathrm{M}_1) + \overline{\alpha}(\mathrm{M}_2)\)。的确，若 \(\mathrm{U}_1\) 与 \(\mathrm{U}_2\) 是 \(\overline{\Phi}\) 的两个开集，分别包含 \(\mathbf{M}_1\) 与 \(\mathbf{M}_2\)，并且使 \(\overline{\alpha} (\mathrm{U}_1)\leqslant \overline{\alpha} (\mathrm{M}_1) + \varepsilon\) 与 \(\overline{\alpha} (\mathrm{U}_2)\leqslant \overline{\alpha} (\mathrm{M}_2) + \varepsilon\)，则 \(\mathrm{U}_1\cup \mathrm{U}_2\) 属于 \(\overline{\Phi}\)，包含 \(\mathbf{M}_1\cup \mathbf{M}_2\)，从而

\[
\overline {{\alpha}} \left(\mathrm{M} _ {1} \cup \mathrm{M} _ {2}\right) \leqslant \overline {{\alpha}} \left(\mathrm{U} _ {1} \cup \mathrm{U} _ {2}\right) \leqslant \overline {{\alpha}} \left(\mathrm{U} _ {1}\right) + \overline {{\alpha}} \left(\mathrm{U} _ {2}\right) \leqslant \overline {{\alpha}} \left(\mathrm{M} _ {1}\right) + \overline {{\alpha}} \left(\mathrm{M} _ {2}\right) + 2 \varepsilon ,
\]

由此即得我们的断言。

其次，我们来证明：若 \(K\) 是紧集而 \(U\) 是 \(\overline{\Phi}\) 的使 \(K \subset U\) 的开集，则 \(\overline{\alpha}(U - K) = \overline{\alpha}(U) - \overline{\alpha}(K)\)。由前所述，我们有 \(\overline{\alpha}(U) \leqslant \overline{\alpha}(K) + \overline{\alpha}(U - K)\)。另一方面，对每个紧集 \(K' \subset U - K\)，

\[
\overline {{{{\alpha}}}} (\mathrm{K} \cup \mathrm{K} ^ {\prime}) = \overline {{{{\alpha}}}} (\mathrm{K}) + \overline {{{{\alpha}}}} (\mathrm{K} ^ {\prime}) \leqslant \overline {{{{\alpha}}}} (\mathrm{U});
\]

由于 U - K 是开集且属于 \(\overline{\Phi}\)，\(\overline{\alpha}(U - K)\) 是诸 \(\overline{\alpha}(K')\) 的上确界，这表明 \(\overline{\alpha}(K) + \overline{\alpha}(U - K) \leqslant \overline{\alpha}(U)\)。

因此，\(\overline{\Phi}\) 的定义现在可以表述如下：为使 \(M \in \overline{\Phi}\)，其充要条件是，对每个 \(\varepsilon > 0\)，存在一个紧集 K 与一个开集 \(U \in \overline{\Phi}\)，使得 \(K \subset M \subset U\) 且 \(\overline{\alpha}(U - K) \leqslant \varepsilon\)。

现在我们能够证明 \(\overline{\Phi}\) 是一个集环，而 \(\overline{\alpha}\) 是 \(\overline{\Phi}\) 上的一个加性集函数。我们首先证明：若 M 与 N 属于 \(\overline{\Phi}\)，则 M \(\cap\) CN 与 M \(\cup\) N 亦然。按假设，对每个 \(\varepsilon > 0\)，存在两个紧集 K,K' 与 \(\overline{\Phi}\) 的两个开集 U,U'，使得

\[
\mathrm{K} \subset \mathrm{M} \subset \mathrm{U}, \quad \mathrm{K} ^ {\prime} \subset \mathrm{N} \subset \mathrm{U} ^ {\prime}, \quad \overline {{\alpha}} (\mathrm{U} - \mathrm{K}) \leqslant \varepsilon , \quad \overline {{\alpha}} (\mathrm{U} ^ {\prime} - \mathrm{K} ^ {\prime}) \leqslant \varepsilon .
\]

集合 \(K'' = K \cap C U'\) 是紧的，集合 \(U'' = U \cap C K'\) 是开集且属于 \(\overline{\Phi}\)，并且 \(K'' \subset M \cap C N \subset U''\)；另一方面，\(U'' - K''\) 含于 \(U \cap C K\) 与 \(U' \cap C K'\) 的并，于是 \(\overline{\alpha}(U'' - K'') \leqslant 2\varepsilon\)，这就证明了 \(M \cap C N \in \overline{\Phi}\)。类似地，\(U_1 = U \cup U'\) 是开集且属于 \(\overline{\Phi}\)，\(K_1 = K \cup K'\) 是紧的，且 \(K_1 \subset M \cup N \subset U_1\)；另一方面，\(U_1 - K_1\) 含于 U - K 与 \(U' - K'\) 的并，于是同样有 \(\overline{\alpha}(U_1 - K_1) \leqslant 2\varepsilon\)，从而 \(M \cup N\) 属于 \(\overline{\Phi}\)。最后，若 M 与 N 不相交，则

\[
\overline {{{{\alpha}}}} \left(\mathrm{K} _ {1}\right) = \overline {{{{\alpha}}}} (\mathrm{K}) + \overline {{{{\alpha}}}} \left(\mathrm{K} ^ {\prime}\right) \geqslant \overline {{{{\alpha}}}} (\mathrm{M}) + \overline {{{{\alpha}}}} (\mathrm{N}) - 2 \varepsilon ,
\]

因而 \(\overline{\alpha} (\mathbf{M}\cup \mathbf{N})\geqslant \overline{\alpha} (\mathbf{M}) + \overline{\alpha} (\mathbf{N}) - 2\varepsilon\)；由于 \(\varepsilon\) 是任意的，我们有 \(\overline{\alpha} (\mathbf{M}\cup \mathbf{N}) = \overline{\alpha} (\mathbf{M}) + \overline{\alpha} (\mathbf{N})\)

\(6^{\circ}\) 测度 \(\mu\) 的存在性。

由 No.~9 的命题 18，存在且仅存在一个正线性形式 \(\beta\)，它定义在向量空间 \(\mathcal{E}(\overline{\Phi})\)（\(\overline{\Phi}\)-阶梯函数所成的）上，使得 \(\beta(\varphi_{\mathrm{M}}) = \overline{\alpha}(\mathrm{M})\) 对所有 \(\mathrm{M} \in \overline{\Phi}\) 成立。对每个紧子集 \(\mathrm{K}\)（\(\mathrm{X}\) 中的），我们用 \(\mathcal{G}(\mathrm{K})\) 表示 \(\mathcal{E}(\overline{\Phi})\) 中支集含于 \(\mathrm{K}\) 的函数的一致极限所成的空间。由于 \(\beta\) 是正的，\(|\beta(f)| \leqslant \overline{\alpha}(\mathrm{K}) \cdot \|f\|\) 对每个函数 \(f \in \mathcal{E}(\overline{\Phi})\)（支集含于 \(\mathrm{K}\) 者）成立；\(\beta\) 在这些函数所成的空间上的限制是关于一致收敛拓扑的一个连续线性形式；因此它可以延拓为一个正连续线性形式 \(\overline{\beta}_{K}\)（\(\mathcal{G}(K)\) 上的）。此外，若 K 与 \(K_{1}\) 是使 \(K \subset K_{1}\) 的两个紧集，则 \(\overline{\beta}_{K_{1}}\) 在 \(\mathcal{G}(K)\) 上的限制与 \(\overline{\beta}_{K}\) 相同，因此存在一个正线性形式 \(\overline{\beta}\)（在诸 \(\mathcal{G}(K)\) 的并 G 上），它延拓每一个形式 \(\overline{\beta}_{K}\)。

现在，由于每个紧集都属于 \(\overline{\Phi}\)，具有紧支集的连续实值函数所成的空间 \(\mathcal{K}\) 是 \(\mathcal{G}\) 的一个子空间（No.~10, 命题 19）；因此，\(\mathcal{K}\) 上的正线性形式 \(\overline{\beta}\) 的限制是一个正测度 \(\mu\)。我们来证明，对每个紧集 \(K\)，有 \(\mu(K) = \overline{\alpha}(K)\)。对每个 \(\varepsilon > 0\)，存在一个开集 \(U \in \overline{\Phi}\)，使得 \(K \subset U\)，\(\mu(U) \leqslant \mu(K) + \varepsilon\) 且 \(\overline{\alpha}(U) \leqslant \overline{\alpha}(K) + \varepsilon\)。设 \(f\) 是 \(X\) 到 {[}0,1{]} 中的一个连续映射，其支集含于 \(U\)，且 \(f(x) = 1\) 在 \(K\) 上成立（Ch. III, §1, No.~2, 引理 1）。那么 \(\mu(K) \leqslant \mu(f) \leqslant \mu(U) \leqslant \mu(K) + \varepsilon\)，而另一方面，

\[
\overline {{\alpha}} (\mathrm{K}) = \beta (\varphi_ {\mathrm{K}}) \leqslant \overline {{\beta}} (f) \leqslant \beta (\varphi_ {\mathrm{U}}) = \overline {{\alpha}} (\mathrm{U}) \leqslant \overline {{\alpha}} (\mathrm{K}) + \varepsilon ;
\]

由于 \(\mu(f) = \overline{\beta}(f)\)，我们看出 \(|\mu(\mathrm{K}) - \overline{\alpha}(\mathrm{K})| \leqslant \varepsilon\)，又因 \(\varepsilon\) 是任意的，故 \(\mu(\mathrm{K}) = \overline{\alpha}(\mathrm{K})\)。

属于 \(\overline{\Phi}\) 的开集的上述刻画，与 No.~6 的定理 4 的推论 4 相结合，便表明属于 \(\overline{\Phi}\) 的开集恰恰就是 \(\mu\)-可积的开集，并且对这样的集合 U，有 \(\mu(\mathrm{U}) = \overline{\alpha}(\mathrm{U})\)。No.~6 的定理 4 与 \(\overline{\Phi}\) 的集合在 \(5^{\circ}\) 中给出的刻画进而表明，\(\mu\)-可积集就是 \(\overline{\Phi}\) 中的集合，并且对这样的集合 M，有 \(\mu(\mathrm{M}) = \overline{\alpha}(\mathrm{M})\)。最后，对每个开集 U，\(\mu^{*}(\mathrm{U}) = \sup_{\mathrm{M} \in \Phi, \mathrm{M} \subset \mathrm{U}} \alpha(\mathrm{M})\) 这一事实由 \((\mathrm{PC}_{\mathrm{II}})\) 与 No.~6 的定理 4 的推论 4 立即得出。

这样，定理 5 就完全证毕。

推论. --- 设 X 是一个具有可数基的局部紧空间，\(\Psi\) 是 X 的 Borel 集所成之集，\(\beta\) 是 \(\Psi\) 到 \([0,+\infty]\) 中的一个映射，满足下列条件：

\begin{enumerate}
\def\labelenumi{(\roman{enumi})}
\item
  若 \((\mathrm{B}_1, \mathrm{B}_2, \ldots)\) 是 X 的两两不相交的 Borel 集的一个序列，则 \(\beta(\mathrm{B}_1 \cup \mathrm{B}_2 \cup \ldots) = \beta(\mathrm{B}_1) + \beta(\mathrm{B}_2) + \ldots\)。
\item
  若 B 是 X 的一个紧子集，则 \(\beta (\mathrm{B}) < +\infty\)
\end{enumerate}

那么，存在且仅存在一个正测度 \(\mu\)（\(X\) 上的），使得 \(\beta(B) = \mu^{*}(B)\) 对所有 \(B \in \Psi\) 成立。

设 \(\Phi\) 是 X 的紧子集所成之集，并设 \(\alpha\) 是 \(\beta\) 在 \(\Phi\) 上的限制。此时条件 \((\mathrm{PC_I})\)、\((\mathrm{PC_{II}})\)、\((\mathrm{PM_I})\)、\((\mathrm{PM_{II}})\)、\((\mathrm{PM_{III}})\) 与 \((\mathrm{PM_{IV}''})\) 均得到满足。设 K 是 X 的一个紧子集，且 \(\varepsilon > 0\)。那么 K 是 X 的相对紧开集的一个递减序列 \((\mathrm{U}_1, \mathrm{U}_2, \ldots)\) 的交（GT, IX, §2, No.~5, 命题 7）。我们有 \(\sum_{n=1}^{\infty} \beta(\mathrm{U}_n - \mathrm{U}_{n+1}) = \beta(\mathrm{U}_1 - \mathrm{K}) < +\infty\)，因此

\[
\beta (\mathrm{U} _ {n}) - \beta (\mathrm{K}) = \beta (\mathrm{U} _ {n} - \mathrm{K}) = \sum_ {p = n} ^ {\infty} \beta (\mathrm{U} _ {p} - \mathrm{U} _ {p + 1})
\]

当 n 趋于 \(\infty\) 时趋于 0。这就证明了条件 \((\mathrm{PM}_{\mathrm{IV}}^{\prime})\) 得到满足。由定理 5，存在 X 上的一个正测度 \(\mu\)，使得对 X 的每个紧子集 K 有 \(\mu(\mathrm{K}) = \alpha(\mathrm{K})\)。由于 X 的每个开集 U 都是紧子集的一个递增序列的并，我们有 \(\mu^{*}(\mathrm{U}) = \beta(\mathrm{U})\)。设 L 是 X 的一个紧子集。由 No.~5 的命题 7，L 的 \(\mu\)-可积子集构成 L 的子集的一个 σ-集环。因此，若 B 是 \(\Psi\) 中含于 L 的一个元素，则 B 是 \(\mu\)-可积的；于是对每个 \(\varepsilon > 0\)，存在 X 中的紧集 K 与开集 U，使得 \(K \subset B \subset U\) 且 \(\mu^{*}(\mathrm{U}) - \mu(\mathrm{K}) \leqslant \varepsilon\)（No.~6, 定理 4）。由于 \(\beta(\mathrm{U}) = \mu^{*}(\mathrm{U})\) 且 \(\beta(\mathrm{K}) = \mu(\mathrm{K})\)，我们看出 \(|\mu^{*}(\mathrm{B}) - \beta(\mathrm{B})| \leqslant 2\varepsilon\)。因此 \(\beta(\mathrm{B}) = \mu^{*}(\mathrm{B})\)。最后，X 的每个 Borel 集C 都是两两不相交的相对紧 Borel 集的一个序列的并，由此 \(\beta(\mathrm{C}) = \mu^{*}(\mathrm{C})\)。\(\mu\) 的唯一性由定理 5 立即得出。

